\chapter{ABSTRACT}

\indent\indent In the era of digital transformation, automatic classification of intangible cultural heritage (ICH) data plays an important role in cultural preservation and promotion. However, traditional CNN-based methods that rely only on image data still face many limitations, especially when dealing with multimodal data and high similarity in costumes, rituals, and festival contexts. To address these issues, I proposes a deep learning architecture that combines a Graph Attention Network on a multimodal graph (MGAT) with FastKAN (Kolmogorov–Arnold Network). The proposed architecture enables joint learning and integration of three data modals, including image, text, and audio, within a shared representation space. Experimental results on a collected ICH dataset show that the proposed method outperforms baseline models, with the DenseNet121 + MGAT + FastKAN configuration achieving an accuracy of 94.35\%. In addition, lightweight models such as MobileNetV3 Small and ShuffleNetV2 1.0x also show improvements of more than 8\%. The use of FastKAN allows the model to learn complex nonlinear relationships more effectively than traditional MLPs. Furthermore, this study develops and deploys a web-based system called ICH Predictor, which supports data lookup and visualization of the ICH classification process, enhancing the practical applicability of the proposed approach.\vspace{0.3cm}

\textbf{Keywords:} Intangible Cultural Heritage, Multimodal Classification, Graph Neural Networks, MGAT, FastKAN, Deep Learning.